\noindent \textit{\textcolor{red!60!black}{
I appreciate using 1936--37 as a placebo. While this helps addressed concerns about downturns driving the results, there are other important macro differences between both periods---such as the collapse of bond issuance, that limited rollover---that may affect exposed firms more in 1933--34 beyond the uncertainty of the abrogation of the gold cause. That is, the placebo helps but only for a subset of the concerns about confounders. To that end, the robustness regressions are more helpful. Though the findings are broadly reassuring, they do point to either the thinness of the data or a lack of robustness. The authors cannot control for industry $\times$ year and $X_{it} \times$ year at the same time, so the fixed effects vary across specifications in Table 6. The paper argues this is done ``To assess robustness to firm-level controls separately from industry-level controls, this specification omits industry-year fixed effects.'' This sounds like a copout. Given the concerns of New Deal policies as confounders, which may have affected firms differently given their characteristics, one may want to control for both. My sense is that they lack power, as even without the industry-year FEs the significance is lost for $1934 \times d$ across several specifications for Columns 2--10. I would like to see the results with industry-year FEs for these columns as well. If significance is lost across the board, then the paper should acknowledge it and make its best case to defend their preferred specifications due to limited power---but readers should know whether the effects are sensitive so they can pass their own judgement.
}}

\bigskip

\noindent We thank the referee for these comments and address each in turn.

\medskip
\noindent \textit{Placebo and the rollover channel.} The referee correctly notes that the 1936--37 placebo addresses recession-driven confounders but not the collapse of bond issuance that may have disproportionately affected gold-exposed firms---which held more bonds---through a rollover channel (the collapse in industrial bond issuance is shown in Figure~7 of the paper). Column 3 of Table 4 speaks to this concern without fully resolving it: it excludes firms with bonds maturing at any point between 1931 and 1934, and the 1933 and 1934 coefficients remain $-0.058$ and $-0.037$, essentially unchanged from the baseline, suggesting that the investment effect is not entirely driven by firms with immediate refinancing needs.

\medskip
\noindent \textit{Industry-year fixed effects in columns 2--10.} We agree that readers should be able to assess this sensitivity directly. Internet Appendix Table~IA.19 re-estimates columns 2--10 of Table 7 (Table 6 in the previous draft) with industry-year fixed effects in place of year fixed effects, so that all ten specifications share the fixed-effect structure of column 1; its columns (1)--(9) correspond to columns 2--10 of Table 7.

\medskip
\noindent The referee's sense that power is the binding constraint is borne out. The paper's main result rests on the baseline of Table 4, whose robustness to industry-year fixed effects column 1 of Table 7 already establishes ($-0.059^{***}$ and $-0.036^{***}$ in 1933 and 1934). Across the nine supplementary specifications, $\tilde{d}$ remains negative and significant at the 5\% level throughout, and $1933 \times \tilde{d}$ at the 5\% level in seven of nine (10\% in all nine). As the referee anticipated, $1934 \times \tilde{d}$ loses significance in seven of nine columns---but the point estimates remain negative and stable ($-0.012$ to $-0.031$, with no sign reversals), so it is the standard errors that widen, not the coefficients that shrink: a power problem rather than a robustness failure.

\medskip
\noindent The precision loss has a mechanical source. Within a 2-digit industry, firms cluster in similar characteristic deciles, so the decile-period dummies and the industry-year cells absorb overlapping variation and standard errors widen; because the 1934 point estimate is smaller than the 1933 one, the widening pushes the 1934 interaction below conventional thresholds first. In the all-controls specification (column (1) of Table~IA.19), the overlap becomes exact collinearity: the four period indicators sum to one, so each characteristic's four interactions sum to the characteristic itself---a firm-level constant that the firm fixed effects already absorb---and once the industry-year cells are absorbed as well, the after-period interaction is dropped for all eight characteristics. That specification is then no longer the fully-controlled model it was designed to be, and its 1933 interaction, like the cash-decile column's, retains significance only at the 10\% level.

\medskip
\noindent Column 1 of Table 7 remains our preferred specification for the New Deal confounder concern: industry-year fixed effects absorb any shock common to an industry in a given year. The characteristic controls address the distinct concern that gold exposure proxies for observables on differential investment trajectories, and they continue to do so under industry-year fixed effects, since $\tilde{d}$ remains significant throughout. Combining the two exercises comes at a power cost; we say so explicitly in Section 5 and provide Table~IA.19 so that readers can judge the sensitivity themselves.
